\section{Theorem B: a 21-parameter certified affine regime}\label{sec:B}
\statusmachine
\subsection{Certification lemmas}
\begin{lemma}[Window repetition]\label{lem:Bwindow}
Fix $k,m$. Suppose for integers $p\ge1,t\ge0$,
$Q_{t+i}=Q_{t+p+i}$ for every $0\le i<k$. Then $Q_n=Q_{n+p}$ for all $n\ge t$.
\end{lemma}
\begin{proof}
The given $k$ equalities form the induction base. For $n\ge t+k$, each row used in the calculation of $A_n$ is $Q_{n-r}$ for $1\le r\le k$, and the corresponding row for $A_{n+p}$ is $Q_{n+p-r}$. These rows agree by induction; therefore their potential-removal sets and the resulting $Q$ rows agree by the exact recurrence in Theorem~\ref{thm:Ainterval}. Since $n\ge t+k\ge k$, both states allow all removal sizes through $k$.
\end{proof}
\begin{lemma}[Starting-position bridge]\label{lem:Bbridge}
Under Lemma~\ref{lem:Bwindow}, the starting-position P/N sequence has period $p$ for all $n\ge t+k$. A proposed $p$-periodic formula is valid for every $n\ge0$ if it matches the exact starting P/N outcome for $0\le n\le t+k+p-1$.
\end{lemma}
\begin{proof}For $n\ge t+k$, each child row $Q_{n-r}$ agrees with $Q_{n+p-r}$, and all removals $1\le r\le k$ are available. Thus $A_n=A_{n+p}$, so the starting outcomes agree. The finite bound includes the entire earlier prefix and one complete period beyond $t+k$; any subsequent index reduces to that verified range.\end{proof}
\begin{theorem}[Exact certified $A=86$ regime]\label{thm:B}
For each integer $58\le k\le78$, the unrestricted starting position of $D(k,1)$ with $n$ counters is losing precisely when
\[\boxed{n\bmod(2k+88)\in\{0,86,k+87\}.}\]
The least positive period of this starting P/N sequence is $2k+88$.
\end{theorem}
\begin{proof}[Finite computation and mathematical induction]
Fix one of the 21 values $k=58,\ldots,78$. The script \texttt{src/verify\_theorems\_ab.py} computes each $Q_n$ by the integer recurrence of Theorem~\ref{thm:Ainterval}. With $p=2k+88$, it finds the exact $t$ shown in Table~\ref{tab:Bcert} and asserts the $k$ row equalities in Lemma~\ref{lem:Bwindow}. It also verifies the proposed start predicate for every $0\le n\le t+k+p-1$. Lemmas~\ref{lem:Bwindow} and~\ref{lem:Bbridge} therefore prove the infinite-$n$ claim for this fixed $k$. The script repeats these finite checks for each of the 21 listed integers. It additionally compares the rows and unrestricted starts with an independent direct Sprague--Grundy implementation over the verified range.

To establish minimality, set $R=\{0,86,k+87\}$ in the cyclic group of order $p$. If a nonzero shift fixes a set of three elements, its translation orbits must each have length three, so $R=\{0,p/3,2p/3\}$. Its smaller nonzero element would therefore be $p/3$. Yet $p/3=(2k+88)/3<86$ for $58\le k\le78$, contradicting $\min(R\setminus\{0\})=86$. Thus the period is least.
\end{proof}
\subsection{Exact reproduced certificates}
\begin{longtable}{rrrrl}
\caption{Theorem B, corrected finite-window certificates. The bridge endpoint is inclusive.}\label{tab:Bcert}\\
\toprule $k$ & $p$ & $t$ & bridge end & independent SG \\
\midrule\endfirsthead\toprule $k$ & $p$ & $t$ & bridge end & independent SG \\
\midrule\endhead
58 & 204 & 58 & 319 & PASS \\
59 & 206 & 59 & 323 & PASS \\
60 & 208 & 59 & 326 & PASS \\
61 & 210 & 59 & 329 & PASS \\
62 & 212 & 59 & 332 & PASS \\
63 & 214 & 58 & 334 & PASS \\
64 & 216 & 59 & 338 & PASS \\
65 & 218 & 0 & 282 & PASS \\
66 & 220 & 0 & 285 & PASS \\
67 & 222 & 0 & 288 & PASS \\
68 & 224 & 0 & 291 & PASS \\
69 & 226 & 0 & 294 & PASS \\
70 & 228 & 0 & 297 & PASS \\
71 & 230 & 0 & 300 & PASS \\
72 & 232 & 0 & 303 & PASS \\
73 & 234 & 0 & 306 & PASS \\
74 & 236 & 0 & 309 & PASS \\
75 & 238 & 0 & 312 & PASS \\
76 & 240 & 0 & 315 & PASS \\
77 & 242 & 0 & 318 & PASS \\
78 & 244 & 0 & 321 & PASS \\
\bottomrule\end{longtable}
\paragraph{Unresolved extension (Working on).} The 21 instances constitute a completed computer-assisted proof, not a uniform symbolic derivation for arbitrary $k$. Explicit P-row formulas explaining the numbers 86 and 88 and the boundary at $k=78$ remain \emph{Working on}. See Theorem C.
